
\begin{center}
\Huge
Polynomier
\end{center}
\section*{Introduktion}
\stepcounter{section}



Vi har tidligere arbejdet med både lineære og konstante funktioner. Disse er begge eksempler på den funktionstype, vi kalder for \textit{polynomier}. Før vi giver en definition på et polynomium, så ser vi på nogle eksempler.

\begin{exa}
	Funktionen $f$ givet ved
	\begin{align*}
		f(x) = 7
	\end{align*}
	er en konstant funktion. Dette kaldes i sammenhæng med polynomier for et \textit{nultegradspolynomium}.
	
	Funktionen $g$ givet ved
	\begin{align*}
		g(x) = 5x -1
	\end{align*}
	er en lineær funktion. Dette kaldes også et \textit{førstegradspolynomium.}
	
	Funktionen $h$ givet ved
	\begin{align*}
		h(x) = 2x^2 -4x + 8
	\end{align*}
	kaldes for et \textit{andengradspolynomium}.
	
	Til slut kaldes funktionen $p$ givet ved
	\begin{align*}
		p(x) = x^7 - 5x^4 + 2x - 10
	\end{align*}
	for et \textit{7.gradspolynomium}.
\end{exa}
Når vi skal afgøre graden af polynomier, så betragter vi den højeste potens af $x$. Konstante funktioners højeste potens af $x$ er 0, da $x^0 = 1$. Vi kunne i princippet skrive $f$ som
\begin{align*}
	f(x) = 10\cdot x^0.
\end{align*}
Tilsvarende er den højeste potens af $x$ for funktionen $g$ ét, da vi kan skrive
\begin{align*}
	g(x) = 5x^1 - 1.
\end{align*}

\begin{defn}
	Et polynomium er en funktion $p$ på formen
	\begin{align*}
		p(x) = a_nx^n+a_{n-1}x^{n-1}+\cdots+a_1 x + a_0,
	\end{align*}
	hvor $a_n \neq 0$ og $n\in \mathbb{N}$. Vi kalder $n$ for graden af $f$. 
\end{defn}
Vi kalder tallene $a_0,a_1,\hdots,a_n$ for et polynomiums koefficienter. På Figur \ref{fig:polys} kan vi se eksempler på et nulte, første, anden og tredjegradspolynomium.

\begin{figure}[H]
\centering
\resizebox{0.45\textwidth}{!}
{
\begin{tikzpicture}
\begin{axis}
[
	axis lines = middle,
	xlabel = $x$, ylabel = $y$,
	xmin = -2, xmax = 3, ymin = -2, ymax = 4
 ]
\addplot[thick, samples = 100, color = olive] {3};
\end{axis}
\end{tikzpicture}
}
\resizebox{0.45\textwidth}{!}
{
\begin{tikzpicture}
\begin{axis}
[
	axis lines = middle,
	xlabel = $x$, ylabel = $y$,
	xmin = -2, xmax = 3, ymin = -2, ymax = 4
 ]
\addplot[thick, samples = 1000, color = olive] {2*x};
\end{axis}
\end{tikzpicture}
}
\resizebox{0.45\textwidth}{!}
{
\begin{tikzpicture}
\begin{axis}
[
	axis lines = middle,
	xlabel = $x$, ylabel = $y$,
	xmin = -2, xmax = 3, ymin = -2, ymax = 4
 ]
\addplot[thick, samples = 1000, color = olive] {-x^2+2*x+1};
\end{axis}
\end{tikzpicture}
}
\resizebox{0.45\textwidth}{!}
{
\begin{tikzpicture}
\begin{axis}
[
	axis lines = middle,
	xlabel = $x$, ylabel = $y$,
	xmin = -2, xmax = 3, ymin = -2, ymax = 4
 ]
\addplot[thick, samples = 1000, color = olive] {x^3-2*x^2+2};
\end{axis}
\end{tikzpicture}
}
\caption{Henholdsvis nulte-, første-, anden- og tredjegradspolynomier.}
\label{fig:polys}
\end{figure}

Vi vil som sagt have et særligt fokus på andengradspolynomier i dette forløb, og derfor betragtes graferne for andengradspolynomier. Disse kaldes også for \textit{parabler}.

For en god ordens skyld giver vi en definition på et andengradspolynomium.
\begin{defn}[Andengradspolynomium]
	En funktion $f$ på formen
	\begin{align*}
		f(x) = ax^2+bx+c,
	\end{align*}
	hvor $a \neq 0$ kaldes for et andengradspolynomium.
\end{defn}
Vi skal kunne bruge koefficienterne i et andengradspolynomium til at afgøre, hvordan grafen for polynomiet ser ud. For koefficienterne gælder der:
\begin{itemize}
	\item[$\cdot$] Tallet $c$ er skæringen med $y$-aksen.
	\item[$\cdot$] Tallet $b$ fortæller os, om parablen er voksende eller aftagende i skæringen
	med $y$-aksen. Hvis $b > 0$, så er parablen voksende, og hvis $b<0$, så er parablen 
	aftagende.
	\item[$\cdot$] Tallet $a$ fortæller os, om parablens arme vender op eller ned. Hvis 
	$a > 0$, så vender armene op ad, og hvis $a < 0$, så vender armene nedad. 
\end{itemize}
På Figur \ref{fig:andenpolys2} kan vi se to parabler. 
\begin{figure}[H]
	\centering
	\resizebox{0.45\textwidth}{!}
	{
	\begin{tikzpicture}
		\begin{axis}[
		axis lines = middle,
		xlabel = $x$, ylabel = $y$,
		xmin = -2, xmax = 6, ymin = -5, ymax = 4,
		ytick = {-2}, xtick = {0},
 		]
			\addplot[thick, samples = 1000, color = olive] {x^2-3*x-2};
			\node[anchor = west] at (axis cs: 4,-2) {$a > 0$};
			\node[anchor = west] at (axis cs: 4,-2.8) {$b < 0$};
			\node[anchor = west] at (axis cs: 4,-3.6) {$c = -2$};
		\end{axis}
	\end{tikzpicture}
	}
	\resizebox{0.45\textwidth}{!}
	{
	\begin{tikzpicture}
		\begin{axis}[
		axis lines = middle,
		xlabel = $x$, ylabel = $y$,
		xmin = -2, xmax = 3, ymin = -2, ymax = 4,
		ytick = {1}, xtick = {0},
 		]
			\addplot[thick, samples = 1000, color = olive] {-x^2+2*x+1};
			\node[anchor = west] at (axis cs: 1,3.8) {$a < 0$};
			\node[anchor = west] at (axis cs: 1,3.3) {$b > 0$};
			\node[anchor = west] at (axis cs: 1,2.8) {$c = 1$};
		\end{axis}
	\end{tikzpicture}
	}
	\caption{To parabler}
	\label{fig:andenpolys2}
\end{figure}

\subsection*{Rødder}
De steder hvor polynomier skærer $x$-aksen kaldes ofte for \textit{rødder}. Vi skal senere se, hvordan vi kan bestemme rødderne for andengradspolynomier, men for nu vil vi holde os til at afgøre, hvor mange rødder et polynomium har. Til dette har vi \textit{diskriminanten}.

\begin{defn}[Rod]
	Lad $p$ være et polynomium. Hvis det for et tal $x_0$ gælder, at
	\begin{align*}
		p(x_0) = 0,
	\end{align*}
	så kaldes $x_0$ for en \textit{rod} for $f$. 
\end{defn}

For andengradspolynomier $p$ givet ved
\begin{align*}
	p(x) = ax^2 + bx + c,
\end{align*}
så er et tal $x_0$ en rod for $p$ hvis og kun hvis $x_0$ er en løsning til andengradsligningen
\begin{align*}
	ax^2 + bx + c = 0.
\end{align*}
Vi kender altså allerede værktøjet til at finde rødder for andengradspolynomier - netop løsningsformlen for andengradsligninger, der som bekendt lyder som følgende men nu formuleret i et polynomiumsperspektiv.
\begin{setn}[Diskriminantformlen]
	Lad et andengradspolynomium $p$ være givet ved
	\begin{align*}
		p(x) = ax^2 + bx + c,
	\end{align*}
	hvor $a \neq 0$. Hvis $d = b^2- 4ac \geq 0$ er rødderne for $p$ givet ved
	\begin{align*}
		x_1 = \frac{-b + \sqrt{d}}{2a}
	\end{align*}
	og
	\begin{align*}
		x_2 = \frac{-b - \sqrt{d}}{2a}.
	\end{align*}
\end{setn}

Parablerne for andengradspolynomier er også med til at forklare, hvorfor andengradsligninger kan have $0$, $1$ og $2$ løsninger. På parablerne i Figurerne \ref{fig:parab1}, \ref{fig:parab2} og \ref{fig:parab3} kan vi netop se, hvorfor andengradsligninger kan have et forskelligt antal løsninger. 
\begin{figure}[H]
	\centering
	\begin{minipage}{0.48\textwidth}
		\centering
	\resizebox{0.95\textwidth}{!}
	{
	\begin{tikzpicture}
	\begin{axis}[
		ticks = none,
		axis lines = center,
		xmin = -4, xmax = 6,
		ymin = -6, ymax = 10,
		xlabel = $x$, ylabel = $y$, 
	]
		\addplot[thick,color = olive, samples = 200, domain = -4:6] {(x-2)*(x+1)};
		\filldraw[color = teal] (axis cs: 2,0) circle (2pt);
		\filldraw[color = teal] (axis cs: -1,0) circle (2pt);
	\end{axis}
	\end{tikzpicture}
	}
	\caption{Andengradspolynomium med to rødder.}
	\label{fig:parab1}
	\end{minipage}
	%%%%
	%%
	%%%%
	\begin{minipage}{0.48\textwidth}
		\centering
	\resizebox{0.95\textwidth}{!}
	{
	\begin{tikzpicture}
	\begin{axis}[
		ticks = none,
		axis lines = center,
		xmin = -4, xmax = 6,
		ymin = -6, ymax = 10,
		xlabel = $x$, ylabel = $y$, 
	]
		\addplot[thick,color = olive, samples = 200, domain = -4:6] {(x-1)^2};
		\filldraw[color = teal] (axis cs: 1,0) circle (2pt);
	\end{axis}
	\end{tikzpicture}
	}
	\caption{Andengradspolynomium med én rod.}
	\label{fig:parab2}
	\end{minipage}
	%%%%
	%%
	%%%%
	\begin{minipage}{0.48\textwidth}
		\centering
	\resizebox{0.95\textwidth}{!}
	{
	\begin{tikzpicture}
	\begin{axis}[
		ticks = none,
		axis lines = center,
		xmin = -4, xmax = 6,
		ymin = -6, ymax = 10,
		xlabel = $x$, ylabel = $y$, 
	]
		\addplot[thick,color = olive, samples = 200, domain = -4:6] {x^2-2*x+3};
	\end{axis}
	\end{tikzpicture}
	}
	\caption{Andengradspolynomium uden rødder.}
	\label{fig:parab3}
	\end{minipage}
\end{figure}

\begin{exa}
	Polynomiet $f$ givet ved
	\begin{align*}
		f(x) = 2x^2-4x+3
	\end{align*}
	har nul rødder, da diskriminanten $d$ er givet ved
	\begin{align*}
		d = (-4)^2 - 4\cdot 2\cdot 3 = 16-24 = -8,
	\end{align*}
	som er mindre end 0.
\end{exa}
\begin{exa}
	Polynomiet $g$ givet ved
	\begin{align*}
		g(x) = -3x^2 + 4x+2
	\end{align*}
	har to rødder, da diskriminanten $d$ er givet ved
	\begin{align*}
		d = 4^2 -4\cdot (-3)\cdot 2 = 16 + 24 = 40,
	\end{align*}
	som er større end 0.
\end{exa}
\begin{exa}
	Polynomiet $h$ givet ved
	\begin{align*}
		h(x) = x^2 - 2x + 1
	\end{align*}
	har netop én rod, da diskriminanten $d$ er givet ved
	\begin{align*}
		d = (-2)^2 - 4\cdot 1 \cdot 1 = 4 - 4 = 0.
	\end{align*}
\end{exa}

\subsection*{Opgave 1}
Bestem graden af følgende polynomier.
\begin{align*}
	&1) \ p(x) = 2x+1   &      &2) \  p(x) = 3x^2+4      \\
	&3) \ p(x) = x^4+2x+10   &      &4) \ p(x) = -10x^5+2x^2+x-0.5      \\
	&5) \ p(x) = 2x+x^2   &      &6) \ p(x) =  x^9     \\
	&7) \ p(x) = -x^4+2x^3-3x^2+10x+2   &      &8) \ p(x) = 5x^2+10x^{12}     \\	
\end{align*}

\subsection*{Opgave 2}
Vi skal afgøre, hvad koefficienterne $a$, $b$ og $c$ betyder for parablen for et andengradspolynomium $p$ givet ved
\begin{align*}
	p(x) = ax^2 + bx + c.
\end{align*}
Følgende skal løses i Maple. 
\begin{enumerate}[label=\roman*)]
	\item Vælg $b$ og $c$ tilfældigt og prøv at tegne grafen for et polynomium med forskellige værdier for $a$. 
	Hvad sker der, hvis $a$ er positiv? Hvad med negativ? Hvis vi gør $a$ stor? Hvis vi gør $a$ lille?
	\item Vælg nu $a$ og $b$ tilfældigt og prøv at ændre på $c$-værdien. Hvad betyder denne for grafen?
	\item Vælg til slut $a$ og $c$ tilfældigt og prøv at ændre på $b$. Kan du gennemskue, hvad $b$ betyder for
	grafen? (Vink: Betragt grafens hældning, der hvor parablen skærer $y$-aksen. Hvad kan vi sige om 
	hældningen, hvis $b$ er positiv? Hvad med negativ?)
\end{enumerate}


\subsection*{Opgave 3}
For parablerne i Figur \ref{fig:andenpolys}
\begin{enumerate}[label=\roman*)]
\item Bestem fortegnet på diskriminanten $d$.
\item Bestem fortegnet på koefficienterne $a$ og $b$. 
\item Bestem $c$.
\end{enumerate}
\begin{figure}[H]
\centering
\resizebox{0.45\textwidth}{!}
{
\begin{tikzpicture}
\begin{axis}
[
	axis lines = middle,
	xlabel = $x$, ylabel = $y$,
	xmin = -2, xmax = 3, ymin = -2, ymax = 4
 ]
\addplot[thick, samples = 1000, color = olive] {x^2-2*x+2};
\end{axis}
\end{tikzpicture}
}
\resizebox{0.45\textwidth}{!}
{
\begin{tikzpicture}
\begin{axis}
[
	axis lines = middle,
	xlabel = $x$, ylabel = $y$,
	xmin = -2, xmax = 3, ymin = -2, ymax = 4
 ]
\addplot[thick, samples = 1000,color = olive] {-2*x^2+4*x-2};
\end{axis}
\end{tikzpicture}
}
\resizebox{0.45\textwidth}{!}
{
\begin{tikzpicture}
\begin{axis}
[
	axis lines = middle,
	xlabel = $x$, ylabel = $y$,
	xmin = -3, xmax = 3, ymin = -2, ymax = 4
 ]
\addplot[thick, samples = 1000, color = olive] {x^2+2*x-1};
\end{axis}
\end{tikzpicture}
}
\resizebox{0.45\textwidth}{!}
{
\begin{tikzpicture}
\begin{axis}
[
	axis lines = middle,
	xlabel = $x$, ylabel = $y$,
	xmin = -2, xmax = 3, ymin = -2, ymax = 4
 ]
\addplot[thick, samples = 1000, color = olive] {-3*x^2-x+2};
\end{axis}
\end{tikzpicture}
}
\caption{Fire parabler}
\label{fig:andenpolys}
\end{figure}


\subsection*{Opgave 4}
Bestem diskriminanten for følgende polynomier og brug den til at afgøre, hvor mange rødder polynomierne har. 
\begin{align*}
&1) \ 2x^2+2x-12   &&2) \ 2x^2-8x+4   \\
&3) \ x^2-4  &&4) \   x^2+3x-4  \\
&5) \ 6x^2-6x-12  &&6) \ 25x^2+100x+100   \\
\end{align*}

\subsection*{Opgave 5}
Skitsér graferne for følgende polynomier og brug din skitse til at afgøre antallet af rødder. Tegn desuden graferne i Maple for at undersøge, om din skitse er korrekt.

\begin{align*}
&1) \  p(x) = x^2-1   &&2) \ p(x) = -10x^2    \\
&3) \  p(x) = -2x^2+2x+2  &&4) \  p(x) = 3x^2-x-3   \\
&5) \ p(x) = -x^2-10    &&6) \  p(x) = x^2+2x+7   \\
\end{align*}

\subsection*{Opgave 6}
Bestem rødderne for følgende andengradspolynomier
\begin{align*}
	&1) \ p(x) = x^2 - 3x - 10  &2) \   p(x) = x^2 + 8x + 12 \\
	&1) \ p(x) = x^2 -6x + 9  &2) \   p(x) = 2x^2-10x-24 \\
\end{align*}